\subsubsection{Manual}\label{trien-khai-audit-manual}

Kiểm toán thủ công được thực hiện trực tiếp trên terminal của host để đối chứng và kiểm tra các cấu hình mà hệ thống tự động không thể quét hoặc yêu cầu đánh giá của quản trị viên.

\textbf{Bảng chi tiết các hoạt động Audit thủ công:}

{\def\LTcaptype{none}
\begin{center}
\captionof{table}{Bảng chi tiết kiểm toán thủ công}\label{tab:chi-tiet-kiem-toan-thu-cong}
\end{center}
\begin{longtable}[]{@{}|
  >{\raggedright\arraybackslash}p{(\linewidth - 8\tabcolsep - 5\arrayrulewidth) * \real{0.0739}}|
  >{\raggedright\arraybackslash}p{(\linewidth - 8\tabcolsep - 5\arrayrulewidth) * \real{0.2995}}|
  >{\raggedright\arraybackslash}p{(\linewidth - 8\tabcolsep - 5\arrayrulewidth) * \real{0.3000}}|
  >{\raggedright\arraybackslash}p{(\linewidth - 8\tabcolsep - 5\arrayrulewidth) * \real{0.3266}}|@{}}
\hline
\begin{minipage}[b]{\linewidth}\raggedright
\textbf{Phần}
\end{minipage} & \begin{minipage}[b]{\linewidth}\raggedright
\textbf{Benchmark}
\end{minipage} & \begin{minipage}[b]{\linewidth}\raggedright
\textbf{Phương pháp thực hiện}
\end{minipage} & \begin{minipage}[b]{\linewidth}\raggedright
\textbf{Hành động thực hiện}
\end{minipage} \\
\hline
\endhead

1 & Phân vùng riêng cho các container & \begin{minipage}[t]{\linewidth}\raggedright
grep \textquotesingle/var/lib/docker\textbackslash s\textquotesingle{} /proc/mounts\\
mountpoint -\/- "\$(docker info -f \textquotesingle\{\{ .DockerRootDir \}\}\textquotesingle)"\strut
\end{minipage} & Chạy lệnh \texttt{mountpoint} để kiểm tra phân vùng riêng biệt của Docker. \\
\hline
& Chỉ cho trusted user kiểm soát Docker Daemon & getent group docker & Kiểm tra các thành viên trong nhóm \texttt{docker} qua tệp \texttt{/etc/group}. \\
\hline
& Cấu hình kiểm tra nhật ký cho Docker Daemon & auditctl -l \textbar{} grep /usr/bin/dockerd & Xác nhận sự hiện diện của quy tắc giám sát cho daemon dockerd. \\
\hline
& Audit folder /run/containerd & auditctl -l \textbar{} grep /run/containerd & Xác nhận sự hiện diện của quy tắc giám sát cho thư mục /run/containerd. \\
\hline
& Audit folder /var/lib/docker & auditctl -l \textbar{} grep /var/lib/docker & Xác nhận sự hiện diện của quy tắc giám sát cho thư mục /var/lib/docker. \\
\hline
& Audit folder /etc/Docker & auditctl -l \textbar{} grep /etc/docker & Xác nhận sự hiện diện của quy tắc giám sát cho thư mục /etc/docker. \\
\hline
& Audit docker.service & \begin{minipage}[t]{\linewidth}\raggedright
systemctl show -p FragmentPath docker.service\\
auditctl -l \textbar{} grep docker.service\strut
\end{minipage} & Kiểm tra sự hiện diện của quy tắc giám sát cho tệp \texttt{docker.service}. \\
\hline
& Audt containerd.sock & \begin{minipage}[t]{\linewidth}\raggedright
grep \textquotesingle containerd.sock\textquotesingle{} /etc/containerd/config.toml\\
auditctl -l \textbar{} grep containerd.sock\strut
\end{minipage} & Kiểm tra sự hiện diện của quy tắc giám sát cho socket \texttt{containerd.sock}. \\
\hline
& Audit docker.sock & \begin{minipage}[t]{\linewidth}\raggedright
systemctl show -p FragmentPath docker.sock\\
auditctl -l \textbar{} grep docker.sock\strut
\end{minipage} & Kiểm tra sự hiện diện của quy tắc giám sát cho socket \texttt{docker.sock}. \\
\hline
& Audt /etc/default/docker & auditctl -l \textbar{} grep /etc/default/docker & Xác nhận sự hiện diện của quy tắc giám sát cho tệp /etc/default/docker. \\
\hline
& Audit /etc/docker/daemon.json & auditctl -l \textbar{} grep /etc/docker/daemon.json & Xác nhận sự hiện diện của quy tắc giám sát cho tệp daemon.json. \\
\hline
& Audit /etc/containerd/config.toml & auditctl -l \textbar{} grep /etc/containerd/config.toml & Xác nhận sự hiện diện của quy tắc giám sát cho tệp config.toml. \\
\hline
& Aduti /etc/sysconfig/docker & auditctl -l \textbar{} grep /etc/sysconfig/docker & Xác nhận sự hiện diện của quy tắc giám sát cho tệp sysconfig/docker. \\
\hline
& Audit /usr/bin/containerd-shim & auditctl -l \textbar{} grep /usr/bin/containerd-shim & Xác nhận sự hiện diện của quy tắc giám sát cho binary containerd-shim. \\
\hline
& Audit /usr/bin/containerd-shim-runc-v1 & auditctl -l \textbar{} grep /usr/bin/containerd-shim-runc-v1 & Xác nhận sự hiện diện của quy tắc giám sát cho containerd-shim-runc-v1. \\
\hline
& Audit /usr/bin/containerd-shim-runc-v2 & auditctl -l \textbar{} grep /usr/bin/containerd-shim-runc-v2 & Xác nhận sự hiện diện của quy tắc giám sát cho containerd-shim-runc-v2. \\
\hline
& Audit /usr/bin/runc & auditctl -l \textbar{} grep /usr/bin/runc & Xác nhận sự hiện diện của quy tắc giám sát cho container runtime runc. \\
\hline
& Đảm báo bảo mật cho container & Kiểm tra thủ công theo security benchmark của hệ thống (CIS, STIG...). Hỏi system admin về tiêu chuẩn áp dụng và xác minh các biện pháp đang thực thi. & Thực hiện kiểm tra các cấu hình hardening trên hệ điều hành host. \\
\hline
& Đảm bảo phiên bản mới cho Docker kt & docker version & Đối chiếu phiên bản hiện tại của Docker với bản phát hành mới nhất. \\
\hline

2 & Chạy docker daemon tránh quyền root & ps -fe \textbar{} grep \textquotesingle dockerd\textquotesingle{} & Kiểm tra tiến trình \texttt{dockerd} để xác định Docker daemon đang chạy với quyền nào. Nếu daemon chạy với quyền root thì cần đánh giá rủi ro hoặc cân nhắc rootless mode nếu hệ thống hỗ trợ. \\
\hline
& Network traffic được đặt hạn chế giữa các container ở default bridge & docker network ls -\/-quiet \textbar{} xargs docker network inspect -\/-format \textquotesingle\{\{ .Name \}\}: \{\{ .Options \}\}\textquotesingle{} & Kiểm tra cấu hình network mặc định của Docker, đặc biệt là khả năng giao tiếp giữa các container trên default bridge. Nếu traffic giữa container không được kiểm soát thì cần cấu hình hạn chế bằng tùy chọn phù hợp. \\
\hline
& Logging level set ở ``info'' & grep "log-level" /etc/docker/daemon.json & Kiểm tra file \texttt{daemon.json} để xác nhận Docker daemon đang dùng mức log \texttt{info}. Nếu đang dùng \texttt{debug} trong môi trường production thì cần điều chỉnh lại. \\
\hline
& Docker được quyền thay đổi iptables & grep -- "-\/-iptables" /etc/docker/daemon.json & Kiểm tra Docker daemon có bị vô hiệu hóa quyền quản lý iptables hay không. Nếu không có giải pháp network thay thế thì không nên tắt chức năng iptables của Docker. \\
\hline
& Đảm bảo không sử dụng registry không an toàn & docker info \textbar{} grep -A5 "Insecure Registries" & Kiểm tra danh sách insecure registries trong Docker. Nếu tồn tại registry không an toàn, cần xác minh có được phê duyệt hay không và hạn chế sử dụng trong production. \\
\hline
& Đảm bảo không sử dụng trình điều khiển aufs & docker info -\/-format \textquotesingle Storage Driver: \{\{ .Driver \}\}\textquotesingle{} & Kiểm tra storage driver hiện tại của Docker. Nếu kết quả là \texttt{aufs}, cần xem xét chuyển sang driver khuyến nghị hơn như \texttt{overlay2}. \\
\hline
& Đảm bảo không dùng devicemapper & docker info -\/-format \textquotesingle Storage Driver: \{\{ .Driver \}\}\textquotesingle{} & Kiểm tra Docker có đang dùng \texttt{devicemapper} hay không. Nếu dùng trong môi trường production không phù hợp, cần đánh giá và chuyển sang storage driver an toàn hơn. \\
\hline
& Đảm bảo xác thực TLS cho Docker Daemon & grep "-\/-tls" /etc/docker/daemon.json & Kiểm tra Docker daemon có cấu hình TLS khi mở API qua TCP hay không. Nếu Docker API lắng nghe từ xa mà không có TLS thì cần cấu hình chứng chỉ và bật xác thực TLS. \\
\hline
& Đảm bảo giới hạn ulimit mặc định & \begin{minipage}[t]{\linewidth}\raggedright
ps -ef \textbar{} grep dockerd\\
grep "default-ulimit" /etc/docker/daemon.json\strut
\end{minipage} & Kiểm tra Docker daemon có cấu hình giới hạn \texttt{ulimit} mặc định hay không. Nếu chưa có giới hạn phù hợp, container có thể sử dụng tài nguyên vượt mức cần thiết. \\
\hline
& Hỗ trợ user namespace & \begin{minipage}[t]{\linewidth}\raggedright
docker info \textbar{} grep -i userns\\
docker info -\/-format \textquotesingle\{\{ .SecurityOptions \}\}\textquotesingle{}\strut
\end{minipage} & Kiểm tra Docker có bật user namespace remapping hay không. Nếu chưa bật, root trong container có thể ánh xạ gần với quyền cao hơn trên host, cần đánh giá theo chính sách bảo mật. \\
\hline
& Cgroup sử dụng được xác nhận & grep "-\/-cgroup-parent" /etc/docker/daemon.json & Kiểm tra Docker daemon có cấu hình \texttt{cgroup-parent} hay không. Nếu hệ thống dùng chính sách quản lý tài nguyên riêng thì cần xác nhận cgroup được cấu hình đúng. \\
\hline
& Đảm bảo device size không thay đổi trừ khi cần thiết & grep "-\/-storage-opt dm.basesize" /etc/docker/daemon.json & Kiểm tra cấu hình base device size. Nếu có thay đổi, cần xác minh lý do vận hành vì thay đổi tùy tiện có thể ảnh hưởng đến bảo mật và hiệu năng. \\
\hline
& Đảm bảo authorization cho Docker client & \begin{minipage}[t]{\linewidth}\raggedright
ps -ef \textbar{} grep dockerd\\
grep "authorization-plugins" /etc/docker/daemon.json\strut
\end{minipage} & Kiểm tra Docker daemon có sử dụng authorization plugin hay không. Nếu môi trường có nhiều người dùng hoặc yêu cầu phân quyền truy cập Docker API, cần cấu hình plugin phù hợp. \\
\hline
& Đảm bảo container bị hạn chế không cho quyền mới & \begin{minipage}[t]{\linewidth}\raggedright
ps -ef \textbar{} grep dockerd\\
grep "no-new-privileges" /etc/docker/daemon.json\strut
\end{minipage} & Kiểm tra cấu hình \texttt{no-new-privileges}. Nếu chưa bật, container có thể có nguy cơ leo thang đặc quyền thông qua các cơ chế như setuid hoặc setgid. \\
\hline
& Đảm bảo container có live restore & \begin{minipage}[t]{\linewidth}\raggedright
docker info -\/-format \textquotesingle\{\{ .LiveRestoreEnabled \}\}\textquotesingle{}\\
grep "live-restore" /etc/docker/daemon.json\strut
\end{minipage} & Kiểm tra tính năng \texttt{live-restore}. Nếu chưa bật, container có thể bị ảnh hưởng khi Docker daemon restart, làm giảm tính sẵn sàng của dịch vụ. \\
\hline
& Đảm bảo có log cho đăng nhập gần và từ xa & \begin{minipage}[t]{\linewidth}\raggedright
docker info -\/-format \textquotesingle\{\{ .LoggingDriver \}\}\textquotesingle{}\\
grep "-\/-log-driver" /etc/docker/daemon.json\strut
\end{minipage} & Kiểm tra logging driver của Docker. Nếu hệ thống yêu cầu giám sát tập trung, cần cấu hình logging driver phù hợp để gửi log về hệ thống quản lý log. \\
\hline
& Đảm bảo Userland Proxy đã tắt & \begin{minipage}[t]{\linewidth}\raggedright
ps -ef \textbar{} grep dockerd\\
grep "userland-proxy" /etc/docker/daemon.json\strut
\end{minipage} & Kiểm tra trạng thái \texttt{userland-proxy}. Nếu không cần thiết, nên tắt để giảm bề mặt tấn công và phụ thuộc vào cơ chế NAT của kernel. \\
\hline
& Đảm bảo daemon-wide custom seccomp profile được thiết đặt & docker info -\/-format \textquotesingle\{\{ .SecurityOptions \}\}\textquotesingle{} & Kiểm tra Docker có sử dụng seccomp profile hay không. Nếu seccomp bị vô hiệu hóa hoặc không có profile phù hợp, cần đánh giá lại chính sách hạn chế system call. \\
\hline
& Đảm bảo các tính năng thử nghiệm không được triển khai trong môi trường sản xuất & docker version -\/-format \textquotesingle\{\{ .Server.Experimental \}\}\textquotesingle{} & Kiểm tra Docker daemon có bật experimental features hay không. Nếu kết quả là \texttt{true} trong môi trường production, cần tắt để tránh rủi ro từ tính năng chưa ổn định. \\
\hline
3 & Đảm bảo quyền root:root cho docker.service & \begin{minipage}[t]{\linewidth}\raggedright
systemctl show -p FragmentPath docker.service\\
stat -c \%U:\%G \textless FragmentPath\textgreater{} \textbar{} grep -v root:root\strut
\end{minipage} & Kiểm tra owner và group của file \texttt{docker.service}. File này phải thuộc sở hữu \texttt{root:root} để tránh người dùng không có quyền chỉnh sửa cấu hình khởi động Docker daemon. \\
\hline
& Đảm bảo docker.service cho quyền hợp lý & \begin{minipage}[t]{\linewidth}\raggedright
systemctl show -p FragmentPath docker.service\\
stat -c \%a \textless FragmentPath\textgreater{}\strut
\end{minipage} & Kiểm tra permission của file \texttt{docker.service}. Quyền truy cập phải là \texttt{644} hoặc hạn chế hơn để tránh sửa đổi trái phép. \\
\hline
& Đảm bảo quyền sở hữu của docker.socket được đặt thành root:root & \begin{minipage}[t]{\linewidth}\raggedright
systemctl show -p FragmentPath docker.socket\\
stat -c \%U:\%G \textless FragmentPath\textgreater{} \textbar{} grep -v root:root\strut
\end{minipage} & Kiểm tra owner và group của file \texttt{docker.socket}. File socket unit cần thuộc sở hữu \texttt{root:root} để bảo vệ điểm giao tiếp với Docker daemon. \\
\hline
& Đảm bảo quyền truy cập tệp docker.socket được đặt thành 644 hoặc hạn chế & \begin{minipage}[t]{\linewidth}\raggedright
systemctl show -p FragmentPath docker.socket\\
stat -c \%a \textless FragmentPath\textgreater{}\strut
\end{minipage} & Kiểm tra permission của file \texttt{docker.socket}. Quyền truy cập phải đủ hạn chế, không cho phép user không liên quan chỉnh sửa socket unit. \\
\hline
& Đảm bảo quyền sở hữu thư mục /etc/docker được đặt thành root:root & stat -c \%U:\%G /etc/docker \textbar{} grep -v root:root & Kiểm tra owner và group của thư mục \texttt{/etc/docker}. Đây là thư mục chứa cấu hình Docker nên phải thuộc sở hữu \texttt{root:root}. \\
\hline
& Đảm bảo quyền truy cập thư mục /etc/docker được đặt thành 755 hoặc hạn chế & stat -c \%a /etc/docker & Kiểm tra permission của thư mục \texttt{/etc/docker}. Quyền truy cập nên là \texttt{755} hoặc hạn chế hơn để bảo vệ cấu hình daemon và chứng chỉ liên quan. \\
\hline
& Đảm bảo quyền truy cập tệp docker.socket & \begin{minipage}[t]{\linewidth}\raggedright
stat -c \%a /var/run/docker.sock\\
stat -c \%a /usr/lib/systemd/\\
system/docker.socket\strut
\end{minipage} & Kiểm tra permission của Docker socket và file unit. Nếu quyền quá rộng, user không được phép có thể truy cập Docker daemon. \\
\hline
& Đảm bảo quyền root:root cho sở hữu tệp chứng chỉ & stat -c \%U:\%G /etc/docker/certs.d/* \textbar{} grep -v root:root & Kiểm tra owner và group của các file chứng chỉ registry. Các file chứng chỉ cần thuộc sở hữu \texttt{root:root} để tránh bị thay đổi trái phép. \\
\hline
& Đảm bảo quyền truy cập chứng chỉ registry được đặt thành 444 hoặc hạn chế hơn & find /etc/docker/certs.d/ -type f -exec stat -c "\%a \%n" \{\} \textbackslash; & Kiểm tra permission của các file chứng chỉ registry. Certificate nên chỉ cho phép đọc và không cho phép ghi bởi user không có quyền. \\
\hline
& Đảm bảo quyền sở hữu Docker server certificate là root:root và file permission 444 & \begin{minipage}[t]{\linewidth}\raggedright
stat -c \%U:\%G \textless path\_to\_tlscert\textgreater{} \textbar{} grep -v root:root\\
stat -c \%a \textless path\_to\_tlscert\textgreater{}\strut
\end{minipage} & Kiểm tra quyền sở hữu và permission của Docker server certificate. Certificate phải thuộc \texttt{root:root} và có permission đọc hạn chế để tránh bị giả mạo hoặc thay đổi. \\
\hline
& Đảm bảo quyền Docker socket file ownership là root:docker và file permission 660 & \begin{minipage}[t]{\linewidth}\raggedright
stat -c \%U:\%G /var/run/docker.sock \textbar{} grep -v root:docker\\
stat -c \%a /var/run/docker.sock\strut
\end{minipage} & Kiểm tra Docker socket thực tế. Socket nên thuộc sở hữu \texttt{root:docker} và permission \texttt{660}, chỉ cho phép root và nhóm docker truy cập. \\
\hline
& Đảm bảo quyền daemon.json là root:root và file permission 644 & \begin{minipage}[t]{\linewidth}\raggedright
stat -c \%U:\%G /etc/docker/daemon.json \textbar{} grep -v root:root\\
stat -c \%a /etc/docker/daemon.json\strut
\end{minipage} & Kiểm tra quyền sở hữu và permission của file \texttt{daemon.json}. Đây là file cấu hình chính của Docker daemon nên phải được bảo vệ chặt chẽ. \\
\hline
& Đảm bảo quyền /etc/sysconfig/docker là root:root và file permission 644 & \begin{minipage}[t]{\linewidth}\raggedright
stat -c \%U:\%G /etc/sysconfig/docker \textbar{} grep -v root:root\\
stat -c \%a /etc/sysconfig/docker\strut
\end{minipage} & Kiểm tra file \texttt{/etc/sysconfig/docker} nếu tồn tại. File này cần thuộc \texttt{root:root} và có permission \texttt{644} hoặc hạn chế hơn. \\
\hline
& Đảm bảo containerd socket có quyền root:root và file permission 660 & \begin{minipage}[t]{\linewidth}\raggedright
stat -c \%U:\%G /run/containerd/containerd.sock \textbar{} grep -v root:root\\
stat -c \%a /run/containerd/containerd.sock\strut
\end{minipage} & Kiểm tra owner, group và permission của \texttt{containerd.sock}. Socket này là điểm giao tiếp với container runtime nên cần giới hạn quyền truy cập. \\
\hline
4 & Phải có container tạo, chạy trong trusted base images và không có package dư thừa không cần thiết được tải về. & \begin{minipage}[t]{\linewidth}\raggedright
docker images\\
docker history \textless imageName\textgreater{}\\
docker exec \$INSTANCE\_ID rpm -qa \# hoặc dpkg -l\strut
\end{minipage} & Kiểm tra image đang sử dụng có đến từ nguồn tin cậy hay không, đồng thời rà soát package bên trong container để phát hiện các gói không cần thiết. \\
\hline
& Đảm bảo image được scan và phải có security patches & \begin{minipage}[t]{\linewidth}\raggedright
docker ps -\/-quiet\\
\# Dùng tool scan (Trivy, Grype, Clair...):\\
trivy image \textless IMAGE\_NAME\textgreater{}\strut
\end{minipage} & Quét image để phát hiện lỗ hổng bảo mật. Nếu có CVE nghiêm trọng, cần cập nhật base image, rebuild image và triển khai lại container. \\
\hline
& Phải có HEALTHCHECK instruction & docker inspect -\/-format=\textquotesingle\{\{ .Config.Healthcheck \}\}\textquotesingle{} \textless IMAGE\textgreater{} & Kiểm tra image có khai báo \texttt{HEALTHCHECK} hay không. Nếu không có, cần bổ sung HEALTHCHECK trong Dockerfile hoặc cấu hình runtime. \\
\hline
& Phải có update instruction và không sử dụng một mình trong Dockerfiles & \begin{minipage}[t]{\linewidth}\raggedright
docker images\\
docker history \textless Image\_ID\textgreater{}\strut
\end{minipage} & Kiểm tra lịch sử build để bảo đảm lệnh update không đứng riêng lẻ. Nếu phát hiện, cần sửa Dockerfile bằng cách gộp update và install trong cùng một lớp build. \\
\hline
& Không cung cấp quyền Setuid và Getuid & docker export \textless IMAGE\_ID\textgreater{} \textbar{} tar -tv 2\textgreater/dev/null \textbar{} grep -E \textquotesingle\^{}{[}-rwx{]}.*(s\textbar S).*\textbackslash s{[}0-9{]}\textquotesingle{} & Kiểm tra các file có setuid/setgid bit trong image. Nếu có file không cần thiết, cần loại bỏ quyền đặc biệt để giảm rủi ro leo thang đặc quyền. \\
\hline
& Dùng COPY thay vì ADD & \begin{minipage}[t]{\linewidth}\raggedright
docker images\\
docker history \textless IMAGE\_ID\textgreater{}\strut
\end{minipage} & Kiểm tra Dockerfile hoặc lịch sử build. Nếu dùng \texttt{ADD} không cần thiết, nên thay bằng \texttt{COPY}. \\
\hline
& Secrets không được bỏ vào Dockerfiles & \begin{minipage}[t]{\linewidth}\raggedright
docker images\\
docker history \textless IMAGE\_ID\textgreater{}\strut
\end{minipage} & Kiểm tra Dockerfile, image layer và lịch sử build để phát hiện password, token, private key hoặc secrets bị nhúng vào image. \\
\hline
& Các package phải chính thống khi tải về & \begin{minipage}[t]{\linewidth}\raggedright
docker images\\
docker history \textless IMAGE\_ID\textgreater{}\strut
\end{minipage} & Kiểm tra nguồn cài đặt package, repository và GPG key. Chỉ sử dụng package từ nguồn đáng tin cậy hoặc repository nội bộ đã phê duyệt. \\
\hline
& Các signed artifacts phải được kiểm chứng & \begin{minipage}[t]{\linewidth}\raggedright
\# Xác minh chữ ký của artifact trước khi upload:\\
cosign verify \textless IMAGE\textgreater{}\\
\# Hoặc kiểm tra chữ ký GPG cho package\strut
\end{minipage} & Kiểm tra chữ ký của image, package hoặc artifact trước khi đưa vào registry hoặc pipeline. Nếu xác minh thất bại thì không được sử dụng artifact đó. \\
\hline
5 & 5.1: Không bật swarm mode nếu không cần thiết & \begin{minipage}[t]{\linewidth}\raggedright
Chạy docker info để đọc trạng thái Swarm.\strut
\end{minipage} & \begin{minipage}[t]{\linewidth}\raggedright
PASS khi Swarm inactive; nếu active thì cần xác nhận nhu cầu sử dụng hoặc ghi FAIL.\strut
\end{minipage} \\
\hline
 & 5.2: Bật AppArmor Profile nếu môi trường hỗ trợ & \begin{minipage}[t]{\linewidth}\raggedright
Dùng docker inspect để đọc AppArmorProfile của từng container.\strut
\end{minipage} & \begin{minipage}[t]{\linewidth}\raggedright
PASS khi mỗi container có profile, ví dụ docker-default.\strut
\end{minipage} \\
\hline
 & 5.3: Thiết lập SELinux security options nếu áp dụng & \begin{minipage}[t]{\linewidth}\raggedright
Chạy getenforce và docker inspect để đọc SecurityOpt, MountLabel, ProcessLabel.\strut
\end{minipage} & \begin{minipage}[t]{\linewidth}\raggedright
Nếu host dùng AppArmor hoặc SELinux disabled thì ghi INFO; nếu áp dụng SELinux thì phải có label phù hợp.\strut
\end{minipage} \\
\hline
 & 5.4: Hạn chế Linux kernel capabilities & \begin{minipage}[t]{\linewidth}\raggedright
Dùng docker inspect để kiểm tra CapAdd và CapDrop.\strut
\end{minipage} & \begin{minipage}[t]{\linewidth}\raggedright
Xác nhận không cấp capability nguy hiểm và không giữ capability mặc định không cần thiết.\strut
\end{minipage} \\
\hline
 & 5.5: Không sử dụng privileged container & \begin{minipage}[t]{\linewidth}\raggedright
Dùng docker inspect để đọc HostConfig.Privileged.\strut
\end{minipage} & \begin{minipage}[t]{\linewidth}\raggedright
PASS khi mọi container có Privileged=false.\strut
\end{minipage} \\
\hline
 & 5.6: Không mount thư mục hệ thống nhạy cảm của host & \begin{minipage}[t]{\linewidth}\raggedright
Dùng docker inspect để kiểm tra danh sách Mounts.\strut
\end{minipage} & \begin{minipage}[t]{\linewidth}\raggedright
Không mount /, /etc, /proc, /sys, /dev hoặc thư mục host nhạy cảm.\strut
\end{minipage} \\
\hline
 & 5.7: Không chạy sshd trong container & \begin{minipage}[t]{\linewidth}\raggedright
Ưu tiên docker top; nếu dùng docker exec thì phải xử lý lỗi thiếu ps.\strut
\end{minipage} & \begin{minipage}[t]{\linewidth}\raggedright
Không có process sshd; nếu không kiểm tra được thì ghi INFO, không tự động PASS.\strut
\end{minipage} \\
\hline
 & 5.8: Không ánh xạ privileged port không cần thiết & \begin{minipage}[t]{\linewidth}\raggedright
Dùng docker inspect để xem port mapping và đối chiếu port nhỏ hơn 1024.\strut
\end{minipage} & \begin{minipage}[t]{\linewidth}\raggedright
Chỉ giữ privileged port đã được phê duyệt; mục chưa tự động hóa thì ghi Manual/INFO.\strut
\end{minipage} \\
\hline
 & 5.9: Chỉ mở port cần thiết & \begin{minipage}[t]{\linewidth}\raggedright
Kiểm tra Config.ExposedPorts và NetworkSettings.Ports.\strut
\end{minipage} & \begin{minipage}[t]{\linewidth}\raggedright
Chỉ publish port có trong thiết kế và tránh mở rộng không cần thiết.\strut
\end{minipage} \\
\hline
 & 5.10: Không chia sẻ network namespace của host & \begin{minipage}[t]{\linewidth}\raggedright
Dùng docker inspect để đọc NetworkMode.\strut
\end{minipage} & \begin{minipage}[t]{\linewidth}\raggedright
PASS khi NetworkMode không phải host.\strut
\end{minipage} \\
\hline
 & 5.11: Giới hạn memory cho container & \begin{minipage}[t]{\linewidth}\raggedright
Dùng docker inspect để đọc HostConfig.Memory.\strut
\end{minipage} & \begin{minipage}[t]{\linewidth}\raggedright
PASS khi memory limit lớn hơn 0 và phù hợp workload.\strut
\end{minipage} \\
\hline
 & 5.12: Đặt CPU priority phù hợp & \begin{minipage}[t]{\linewidth}\raggedright
Dùng docker inspect để đọc CpuShares hoặc CPU quota.\strut
\end{minipage} & \begin{minipage}[t]{\linewidth}\raggedright
Ghi PASS khi có giá trị CPU rõ ràng; giá trị mặc định cần review.\strut
\end{minipage} \\
\hline
 & 5.13: Root filesystem ở chế độ read-only & \begin{minipage}[t]{\linewidth}\raggedright
Dùng docker inspect để đọc ReadonlyRootfs.\strut
\end{minipage} & \begin{minipage}[t]{\linewidth}\raggedright
PASS khi ReadonlyRootfs=true và vùng cần ghi được tách riêng.\strut
\end{minipage} \\
\hline
 & 5.14: Bind traffic vào interface cụ thể & \begin{minipage}[t]{\linewidth}\raggedright
Kiểm tra HostIp trong NetworkSettings.Ports.\strut
\end{minipage} & \begin{minipage}[t]{\linewidth}\raggedright
PASS khi không bind 0.0.0.0, trừ khi có phê duyệt.\strut
\end{minipage} \\
\hline
 & 5.15: Giới hạn restart policy & \begin{minipage}[t]{\linewidth}\raggedright
Đọc RestartPolicy.Name và MaximumRetryCount.\strut
\end{minipage} & \begin{minipage}[t]{\linewidth}\raggedright
PASS khi on-failure và retry từ 1 đến 5.\strut
\end{minipage} \\
\hline
 & 5.16: Không chia sẻ PID namespace của host & \begin{minipage}[t]{\linewidth}\raggedright
Dùng docker inspect để đọc PidMode.\strut
\end{minipage} & \begin{minipage}[t]{\linewidth}\raggedright
PASS khi PidMode không phải host.\strut
\end{minipage} \\
\hline
 & 5.17: Không chia sẻ IPC namespace của host & \begin{minipage}[t]{\linewidth}\raggedright
Dùng docker inspect để đọc IpcMode.\strut
\end{minipage} & \begin{minipage}[t]{\linewidth}\raggedright
PASS khi IpcMode private hoặc không phải host.\strut
\end{minipage} \\
\hline
 & 5.18: Không expose trực tiếp host device & \begin{minipage}[t]{\linewidth}\raggedright
Dùng docker inspect để đọc Devices.\strut
\end{minipage} & \begin{minipage}[t]{\linewidth}\raggedright
PASS khi Devices rỗng hoặc không có thiết bị host chưa phê duyệt.\strut
\end{minipage} \\
\hline
 & 5.19: Ghi đè ulimit khi cần & \begin{minipage}[t]{\linewidth}\raggedright
Dùng docker inspect để đọc Ulimits.\strut
\end{minipage} & \begin{minipage}[t]{\linewidth}\raggedright
Ghi PASS khi có ulimit phù hợp hoặc daemon default đã được harden.\strut
\end{minipage} \\
\hline
 & 5.20: Không dùng shared mount propagation & \begin{minipage}[t]{\linewidth}\raggedright
Dùng docker inspect để đọc Propagation trong Mounts.\strut
\end{minipage} & \begin{minipage}[t]{\linewidth}\raggedright
PASS khi propagation là private, slave hoặc không shared.\strut
\end{minipage} \\
\hline
 & 5.21: Không chia sẻ UTS namespace của host & \begin{minipage}[t]{\linewidth}\raggedright
Dùng docker inspect để đọc UTSMode.\strut
\end{minipage} & \begin{minipage}[t]{\linewidth}\raggedright
PASS khi UTSMode không phải host.\strut
\end{minipage} \\
\hline
 & 5.22: Không tắt default seccomp profile & \begin{minipage}[t]{\linewidth}\raggedright
Dùng docker inspect để đọc SecurityOpt.\strut
\end{minipage} & \begin{minipage}[t]{\linewidth}\raggedright
PASS khi không có seccomp=unconfined.\strut
\end{minipage} \\
\hline
 & 5.23: Không dùng docker exec với privileged option & \begin{minipage}[t]{\linewidth}\raggedright
Kiểm tra journalctl hoặc audit log cho mẫu exec privileged.\strut
\end{minipage} & \begin{minipage}[t]{\linewidth}\raggedright
PASS chỉ áp dụng trong cửa sổ log được truy vấn và nguồn log đáng tin cậy.\strut
\end{minipage} \\
\hline
 & 5.24: Không dùng docker exec với user root option & \begin{minipage}[t]{\linewidth}\raggedright
Kiểm tra journalctl hoặc audit log cho mẫu exec user root.\strut
\end{minipage} & \begin{minipage}[t]{\linewidth}\raggedright
PASS khi không có root exec trong log; cần logging tập trung để tăng độ tin cậy.\strut
\end{minipage} \\
\hline
 & 5.25: Xác nhận cgroup usage & \begin{minipage}[t]{\linewidth}\raggedright
Dùng docker inspect để đọc CgroupParent.\strut
\end{minipage} & \begin{minipage}[t]{\linewidth}\raggedright
Default hierarchy là PASS; custom cgroup phải có lý do và resource control.\strut
\end{minipage} \\
\hline
 & 5.26: Ngăn container nhận thêm đặc quyền & \begin{minipage}[t]{\linewidth}\raggedright
Dùng docker inspect để tìm no-new-privileges trong SecurityOpt.\strut
\end{minipage} & \begin{minipage}[t]{\linewidth}\raggedright
PASS khi mọi container có no-new-privileges.\strut
\end{minipage} \\
\hline
 & 5.27: Có runtime health check & \begin{minipage}[t]{\linewidth}\raggedright
Dùng docker inspect để đọc State.Health.Status.\strut
\end{minipage} & \begin{minipage}[t]{\linewidth}\raggedright
PASS nên là healthy; starting là INFO, unhealthy là FAIL/WARN.\strut
\end{minipage} \\
\hline
 & 5.28: Ghim phiên bản image, không dùng latest & \begin{minipage}[t]{\linewidth}\raggedright
Dùng docker inspect để đọc Config.Image.\strut
\end{minipage} & \begin{minipage}[t]{\linewidth}\raggedright
PASS khi image có version tag cụ thể hoặc digest sha256.\strut
\end{minipage} \\
\hline
 & 5.29: Sử dụng PIDs cgroup limit & \begin{minipage}[t]{\linewidth}\raggedright
Dùng docker inspect để đọc PidsLimit.\strut
\end{minipage} & \begin{minipage}[t]{\linewidth}\raggedright
PASS khi PidsLimit là số dương phù hợp workload.\strut
\end{minipage} \\
\hline
 & 5.30: Không dùng default bridge docker0 & \begin{minipage}[t]{\linewidth}\raggedright
Dùng docker inspect để đọc NetworkSettings.Networks.\strut
\end{minipage} & \begin{minipage}[t]{\linewidth}\raggedright
PASS khi container dùng user-defined network như secure-net.\strut
\end{minipage} \\
\hline
 & 5.31: Không chia sẻ user namespace của host & \begin{minipage}[t]{\linewidth}\raggedright
Dùng docker inspect để đọc UsernsMode.\strut
\end{minipage} & \begin{minipage}[t]{\linewidth}\raggedright
PASS khi UsernsMode không phải host.\strut
\end{minipage} \\
\hline
 & 5.32: Không mount Docker socket vào container & \begin{minipage}[t]{\linewidth}\raggedright
Dùng docker inspect Mounts và tìm docker.sock.\strut
\end{minipage} & \begin{minipage}[t]{\linewidth}\raggedright
PASS khi không mount /var/run/docker.sock hoặc socket tương đương.\strut
\end{minipage} \\
\hline
6 & Tránh Image sprawl & \begin{minipage}[t]{\linewidth}\raggedright
docker images -\/-quiet \textbar{} xargs docker inspect -\/-format \textquotesingle\{\{ .Id \}\}: Image=\{\{ .Config.Image \}\}\textquotesingle{}\\
docker images\strut
\end{minipage} & Rà soát toàn bộ các images hiện có trên máy chủ để phát hiện và gỡ bỏ các image không sử dụng. \\
\hline
& Tránh container sprawl & \begin{minipage}[t]{\linewidth}\raggedright
docker info -\/-format \textquotesingle\{\{ .Containers \}\}\textquotesingle{}\\
docker info -\/-format \textquotesingle\{\{ .ContainersStopped \}\}\textquotesingle{}\\
docker info -\/-format \textquotesingle\{\{ .ContainersRunning \}\}\textquotesingle{}\strut
\end{minipage} & Rà soát tất cả các container đang chạy và đã dừng để loại bỏ các container dư thừa. \\
\hline
\end{longtable}
}
